\subsection{商空间}

设 X 是道路连通的拓扑空间，带有一个离散群 G 的逆紧（右）作用。令 Y = X/G，并用 \(f: X \to Y\) 表示规范映射。若 \(g \in G\) 且 c: I → X 是 X 中的一条道路，则用 \(g^{*}c\) 表示道路 \(t \mapsto c(t) \cdot g\)，用 \(g^{*}[c]\) 表示它的严格同伦类。

设 \(o\) 是 X 的一点。对每个 \(g \in G\)，令 \(\beta_{g}\) 为 X 中连接 \(o \cdot g\) 与 o 的一条道路的类。对每个 \(g \in G\)，令 \(X^{g}\) 为点 \(x \in X\)（满足 \(x \cdot g = x\)）组成的集合；对 \(X^{g}\) 的每个道路连通分支 j，令 \(a_{j}\) 为 j 的一点，并令 \(\gamma_{j}\) 为 X 中连接 \(a_{j}\) 与 o 的一条道路的类。令 \(\nu: F(G) \to \pi_{1}(Y, f(o))\) 为唯一的群同态，使得 \(\nu(g) = f_{*}(\beta_{g})\)（对 \(g \in G\)）。令 N: \(\pi_{1}(X, o) * F(G) \to \pi_{1}(Y, f(o))\) 为唯一的群同态，它与 \(\pi_{1}(f, o)\) 在 \(\pi_{1}(X, o)\) 上、与 \(\nu\) 在 \(F(G)\) 上一致。

命题 3. --- 假设 X 是可解圈的。于是同态 N 是满射的，其核是 \(\pi_{1}(\mathrm{X},o)*\mathrm{F}(\mathrm{G})\) 中包含元素

\[
r _ {2} (k, v) = [ k ] ^ {- 1} v [ k ] (\beta_ {k} ^ {- 1} k ^ {*} (v) ^ {- 1} \beta_ {k})\tag{\((R_{2})\)}
\]

\[
\text {for } k \in \mathrm{G} \text { and } v \in \pi_ {1} (\mathrm{X}, o);\tag{\((R_{3}^{\prime})\)}
\]

\[
r _ {3} ^ {\prime} (k, j) = [ k ] (\beta_ {k} ^ {- 1} k ^ {*} (\gamma_ {j}) ^ {- 1} \gamma_ {j})
\]

的最小正规子群，其中 \(k\in \mathrm{G}\)，\(j\in \pi_0(\mathrm{X}^k)\)

\[
r _ {3} ^ {\prime \prime} (k, h) = [ k h ] ^ {- 1} [ k ] [ h ] (\beta_ {h} ^ {- 1} h ^ {*} (\beta_ {k} ^ {- 1}) \beta_ {k h})\tag{\((\mathrm{R}_{3}^{\prime \prime})\)}
\]

\[
\text {for } k \text { and } h \in G.
\]

群胚态射 \(\varpi(f)\) 诱导出一个态射

\[
\varpi^ {\prime \prime} (f) \colon \varpi (\mathrm{X}) / \mathrm{G} \to \varpi (\mathrm{Y})
\]

由定理 3（IV, 第~403~页），由于空间 X 假设为可解圈的，它是同构。命题于是由 II, 第~211~页命题 6 得出。

下面三个推论由 II, 第~211~页命题 6 的相应推论直接得出。

推论 1. --- 假设 X 是可解圈的，且群 G 由 X 的诸点的稳定子生成。于是规范态射 \(\pi_{1}(f,o)\colon\pi_{1}(\mathrm{X},o)\to\pi_{1}(\mathrm{Y},f(o))\) 是满射的。特别地，若 X 是道路单连通的，则 Y 亦然。

注. --- 若 X 是可解圈的，则 Y 也是（IV, 第~349~页，命题 8）。由于连通的可解圈空间是道路单连通的当且仅当它是单连通的（IV, 第~344~页，定理 1 的推论 1），我们这样就重新得到 I, 第~137~页的命题 11。

例 1. --- 设 X 是道路连通、可解圈且分离的拓扑空间，a 是 X 的一点。设 n 是整数 \(\geqslant 2\)，Y 是空间 \(X^{n}\) 关于群 \(S_{n}\) 通过因子置换而作用所得的商；用 \(f: X^{n} \to Y\) 表示规范映射；用 \(g: X \to Y\) 表示映射 \(x \mapsto f(x, a, \ldots, a)\)。由该命题得出，对每个 i，同态 \(\pi_{1}(g, a)\)（从 \(\pi_{1}(X, a)\) 到 \(\pi_{1}(Y, g(a))\) 中）是满射的，其核是 \(\pi_{1}(X, a)\) 的导出子群。特别地，群 \(\pi_{1}(Y, g(a))\) 是交换的。

推论 2. --- 假设 X 是可解圈的，且群 G 自由作用于 X。存在唯一的群同态 \(p \colon \pi_1(\mathrm{Y}, f(o)) \to \mathrm{G}\)，其核包含 \(\pi_1(\mathrm{X}, o)\) 的像，且使得 \(p(\mathsf{N}(g)) = g\)（对所有 \(g \in \mathrm{G}\)）。此外，\(\pi_1(\mathrm{X}, o) \to \pi_1(\mathrm{Y}, f(o)) \xrightarrow{p} \mathrm{G}\) 是 G 借助 \(\pi_1(\mathrm{X}, o)\) 的一个扩张。

推论 3. --- 假设 X 是道路单连通的。从 G 到 \(\pi_1(Y, f(o))\)、把 \(g \in G\) 对应为道路类 \(f_*(\beta_g)\) 的映射是满射群同态；其核是由 X 的诸点的稳定子生成的 G 的子群。

